\section{Conclusions: resonance theory as a connection problem}
\label{sec:conclusions}
% BEGIN CURATED SUBJECT INDEX
\index{operator!adjoint and self-adjoint}
\index{unbounded operator}
\index{Hilbert space}
\index{locally convex space}
\index{topology}
\index{measure theory}
\index{spectral measure}
\index{spectrum!point}
\index{resolvent}
\index{direct integral}
% END CURATED SUBJECT INDEX

Quantum resonance is a meeting point of scattering theory, non-self-adjoint
spectral analysis, asymptotics, and time-dependent response.  The calculation
is least ambiguous when it starts from Borel-resummed local branches and their
global connection matrix.  The no-incoming condition selects a zero of an
analytic coefficient; the continued resolvent and $S$ matrix have a pole at
that zero.  Feshbach projection packages the same continuum dependence into an
energy-dependent operator, complex scaling exposes the pole as an $L^2$
eigenvalue, Zel'dovich regularization defines its bilinear normalization, and
a rigged Hilbert space places its residue vector in a continuous dual.  These
are not competing definitions but calculational realizations reached after
the invariant connection datum has been fixed.

Functional analysis makes this dictionary more physical, not less.  The
operator domain specifies the observable, topology specifies the admissible
limit, the spectral measure specifies the continuum, and a locally convex test
space specifies which generalized eigenvectors are continuous.  In the direct-
integral spectral representation, energy is the base measure and open channels
form the fibers; the scattering operator is literally the direct integral of
the on-shell matrices.  \vN\ algebras extend this organization to
factor and superselection decompositions.  Resonances sit at the controlled
interface between that measurable self-adjoint structure and analytic
continuation.

The framework remains useful when the problem is enlarged.  A radial singular
point contributes monodromy that cannot be discarded as a formally higher-order term.  An exceptional point turns pole continuation into a branched
Jordan problem and forces continuum normalization to be regulated.  A
black-hole quasinormal mode is an outgoing resonance on the tortoise line, and
complex scaling treats its pole and continuum response within one spectral
representation.

Long-range forces also show where the elementary picture must change.  A
two-body Coulomb state is asymptotic to a logarithmically distorted Coulomb
wave, not to a free plane wave.  Three charged particles require an atlas of
breakup and cluster asymptotics rather than a universal product state.  In
QED, Gauss' law and the soft cloud can move charged asymptotic states out of
the vacuum Fock representation altogether.  The growing complexity with
particle number is therefore not merely combinatorial: long-range
correlations change the representation and the composite operator that
creates the channel.

The unification does not erase the assumptions of the individual methods.
Dilation analyticity, test-function topology, threshold sheets, Borel
directions, and numerical non-normality all matter.  Stating those assumptions
is what converts a suggestive analogy into a transferable method.  The most
promising next steps are to compare residues and continuum response---not only
pole positions---across complex scaling and exact WKB, to construct uniform
long-range channel asymptotics for charged few-body systems, and to formulate
the corresponding resonance problem between infrared representations in
field theory.
